% GENERATED FILE. Edit canonical lessons, then run npm run generate:books.
% canonical-source-hash: fnv1a64:16356c9849689c8a
% canonical-lessons: IT-C455-calzolaio, IT-C455-risuolare, IT-C455-sarto, IT-C455-accorciare, IT-C455-fabbro

\chapter{Riparare e sistemare}
\label{ch:it-riparare-e-sistemare}
\hlchaptermodality{455}

\begin{chapteropening}
\textbf{By the end of this chapter:} \emph{I can get things mended: the cobbler, resoling, the tailor, taking up trousers and the locksmith.}

Complete the last lesson of chapter 455: I can get things mended: the cobbler, resoling, the tailor, taking up trousers and the locksmith.
\end{chapteropening}

\section[il calzolaio]{il calzolaio — the cobbler, the shoemaker}

\label{lesson:IT-C455-calzolaio}



\begin{quote}
Before the new one: say the Italian for the cake shop, then the Italian for the rotisserie.
\end{quote}

\subsection*{You'll want to know: il calzolaio}
\textbf{il calzolaio} — \textquotedblleft{}the cobbler, the shoemaker\textquotedblright{}.

\begin{quote}
Il calzolaio ripara le mie scarpe vecchie.
\end{quote}

\begin{quote}
\emph{The cobbler repairs my old shoes.}
\end{quote}

\emph{Le calze} — socks: both go back to Latin \emph{calceus}, shoe.

\begin{cousinweb}
\emph{calzolaio} is built on \emph{calzare} 'to put on shoes', which is inherited from Latin \emph{calceāre} (EtymDB), from \emph{calceus} 'shoe' (Lewis \& Short).
\end{cousinweb}

\begin{grammarlens}[title={What you've built}]
One more word for telling, explaining and describing at B1.
\end{grammarlens}

\subsection*{Your turn}
\begin{itemize}
\raggedright
  \item \emph{Say it:} \emph{il calzolaio}
  \item \emph{Say it:} \emph{il calzolaio}, once more
  \item \emph{Say it:} a sentence of your own with \emph{il calzolaio}
  \item \emph{Recall:} say the Italian for the cake shop, then the Italian for the rotisserie, then say \emph{il calzolaio} again
\end{itemize}

\begin{tcolorbox}[breakable,colback=teal!4,colframe=teal!35!black,title={Before you move on}]
What does il calzolaio mean? (\textquotedblleft{}The cobbler, the shoemaker\textquotedblright{}.) Say it once more.
\end{tcolorbox}
\section[risuolare]{risuolare — to resole}

\label{lesson:IT-C455-risuolare}



\begin{quote}
Before the new one: say the Italian for the rotisserie, then the Italian for the cobbler.
\end{quote}

\subsection*{You'll want to know: risuolare}
\textbf{risuolare} — \textquotedblleft{}to resole\textquotedblright{}.

\begin{quote}
Faccio risuolare gli stivali.
\end{quote}

\begin{quote}
\emph{I am having my boots resoled.}
\end{quote}

\emph{La suola} — the sole.

\begin{cousinweb}
\emph{risuolare} is built from \emph{ri-} 'again' and \emph{suola} 'sole', which is inherited from Latin \emph{solea} 'sandal, sole', built on \emph{solum} 'ground' (Lewis \& Short).
\end{cousinweb}

\begin{grammarlens}[title={What you've built}]
One more word for telling, explaining and describing at B1.
\end{grammarlens}

\subsection*{Your turn}
\begin{itemize}
\raggedright
  \item \emph{Say it:} \emph{risuolare}
  \item \emph{Say it:} \emph{risuolare}, once more
  \item \emph{Say it:} a sentence of your own with \emph{risuolare}
  \item \emph{Recall:} say the Italian for the rotisserie, then the Italian for the cobbler, then say \emph{risuolare} again
\end{itemize}

\begin{tcolorbox}[breakable,colback=teal!4,colframe=teal!35!black,title={Before you move on}]
What does risuolare mean? (\textquotedblleft{}To resole\textquotedblright{}.) Say it once more.
\end{tcolorbox}
\section[il sarto]{il sarto — the tailor}

\label{lesson:IT-C455-sarto}



\begin{quote}
Before the new one: say the Italian for the cobbler, then the Italian for to resole.
\end{quote}

\subsection*{You'll want to know: il sarto}
\textbf{il sarto} — \textquotedblleft{}the tailor\textquotedblright{}.

\begin{quote}
Il sarto mi fa un vestito per il matrimonio.
\end{quote}

\begin{quote}
\emph{The tailor is making me a suit for the wedding.}
\end{quote}

A woman is \emph{la sarta}.

\begin{cousinweb}
\emph{sarto} is inherited from Late Latin \emph{sartor} 'tailor' (EtymDB), from \emph{sarcīre} 'to patch, to mend'.
\end{cousinweb}

\begin{grammarlens}[title={What you've built}]
One more word for telling, explaining and describing at B1.
\end{grammarlens}

\subsection*{Your turn}
\begin{itemize}
\raggedright
  \item \emph{Say it:} \emph{il sarto}
  \item \emph{Say it:} \emph{il sarto}, once more
  \item \emph{Say it:} a sentence of your own with \emph{il sarto}
  \item \emph{Recall:} say the Italian for the cobbler, then the Italian for to resole, then say \emph{il sarto} again
\end{itemize}

\begin{tcolorbox}[breakable,colback=teal!4,colframe=teal!35!black,title={Before you move on}]
What does il sarto mean? (\textquotedblleft{}The tailor\textquotedblright{}.) Say it once more.
\end{tcolorbox}
\section[accorciare]{accorciare — to shorten, to take up}

\label{lesson:IT-C455-accorciare}



\begin{quote}
Before the new one: say the Italian for to resole, then the Italian for the tailor.
\end{quote}

\subsection*{You'll want to know: accorciare}
\textbf{accorciare} — \textquotedblleft{}to shorten, to take up\textquotedblright{}.

\begin{quote}
Mi può accorciare questi pantaloni?
\end{quote}

\begin{quote}
\emph{Can you take up these trousers for me?}
\end{quote}

\emph{Accorciare la strada} — to take a short cut.

\begin{cousinweb}
\emph{accorciare} is built from \emph{a-} and \emph{corto} 'short', which is inherited from Latin \emph{curtus} 'cut short' (Lewis \& Short).
\end{cousinweb}

\begin{grammarlens}[title={What you've built}]
One more word for telling, explaining and describing at B1.
\end{grammarlens}

\subsection*{Your turn}
\begin{itemize}
\raggedright
  \item \emph{Say it:} \emph{accorciare}
  \item \emph{Say it:} \emph{accorciare}, once more
  \item \emph{Say it:} a sentence of your own with \emph{accorciare}
  \item \emph{Recall:} say the Italian for to resole, then the Italian for the tailor, then say \emph{accorciare} again
\end{itemize}

\begin{tcolorbox}[breakable,colback=teal!4,colframe=teal!35!black,title={Before you move on}]
What does accorciare mean? (\textquotedblleft{}To shorten, to take up\textquotedblright{}.) Say it once more.
\end{tcolorbox}
\section[il fabbro]{il fabbro — the locksmith; the blacksmith}

\label{lesson:IT-C455-fabbro}



\begin{quote}
Before the new one: say the Italian for the tailor, then the Italian for to shorten.
\end{quote}

\subsection*{You'll want to know: il fabbro}
\textbf{il fabbro} — \textquotedblleft{}the locksmith; the blacksmith\textquotedblright{}.

\begin{quote}
Ho perso le chiavi e ho chiamato il fabbro.
\end{quote}

\begin{quote}
\emph{I lost my keys and called the locksmith.}
\end{quote}

Today mostly the locksmith; once the blacksmith.

\begin{cousinweb}
\emph{fabbro} is inherited from Latin \emph{faber, fabrum} 'craftsman, smith' (EtymDB).
\end{cousinweb}

\begin{grammarlens}[title={What you've built}]
One more word for telling, explaining and describing at B1.
\end{grammarlens}

\subsection*{Your turn}
\begin{itemize}
\raggedright
  \item \emph{Say it:} \emph{il fabbro}
  \item \emph{Say it:} \emph{il fabbro}, once more
  \item \emph{Say it:} a sentence of your own with \emph{il fabbro}
  \item \emph{Recall:} say the Italian for the tailor, then the Italian for to shorten, then say \emph{il fabbro} again
\end{itemize}

\begin{tcolorbox}[breakable,colback=teal!4,colframe=teal!35!black,title={Before you move on}]
What does il fabbro mean? (\textquotedblleft{}The locksmith; the blacksmith\textquotedblright{}.) Say it once more.
\end{tcolorbox}
